\documentclass[12pt,a4paper]{article}

\usepackage[utf8]{inputenc}
\usepackage{lmodern}
\usepackage{amsmath, amssymb}
\usepackage{geometry}
\usepackage{enumitem}
\usepackage{fancyhdr}
\usepackage{xcolor}
\usepackage{graphicx}
\usepackage{hyperref}

% Page geometry configuration
\geometry{top=3cm, bottom=2.5cm, left=2.5cm, right=2.5cm}

% Header and Footer configuration
\pagestyle{fancy}
\fancyhf{}
\fancyhead[L]{\textbf{Olavo Barros (102014)}}
\fancyhead[R]{\textbf{Proposta de Projeto}}
\fancyfoot[C]{\thepage}
\renewcommand{\headrulewidth}{0.4pt}
\renewcommand{\footrulewidth}{0.4pt}
\setlength{\headheight}{15pt}

\begin{document}

% Custom Title Block
\begin{center}
    {\LARGE \textbf{Proposta de Projeto}}\\[0.5em]
    {\large \textbf{Dado Tabular Temporal para CNN 2D em FPGA}}\\[1em]
    \textcolor{gray}{\rule{\textwidth}{1pt}}\\[1em]
    \textbf{Aluno:} Olavo Barros \quad \textbf{Matrícula:} 102014\\[0.5em]
    \textbf{Disciplina:} Disciplina de CNN \\[0.5em]
    DPI - Universidade Federal de Viçosa (UFV)
\end{center}
\vspace{1.5em}

\section*{Tema do Projeto}
Aceleração \textit{Hardware-Aware} de Redes Neurais Convolucionais para Classificação de Séries Temporais via Transformação 2D.

\section*{Escopo do Projeto}
O escopo concentra-se no mapeamento de dependências temporais em estruturas geométricas, superando a limitação sequencial associada a dados temporais puros. O projeto é dividido nas seguintes fases principais:

\begin{itemize}[leftmargin=*]
    \item \textbf{Fase 1: Transformação de Dados (1D para 2D):} Aplicação do \textit{Short-Time Fourier Transform} (STFT) para converter janelas do sinal temporal contínuo em matrizes bidimensionais (espectrogramas). Essa abordagem condensa a complexidade temporal em um arranjo espacial estático.
    \item \textbf{Fase 2: Arquitetura da CNN 2D e \textit{TinyML}:} Construção de uma CNN altamente otimizada, desenhada com foco em limitações físicas de hardware. Para acomodar as características dos espectrogramas, serão utilizados kernels assimétricos (ex., 17x9), minimizando a necessidade de parâmetros e ciclos de acesso à memória.
\end{itemize}

\section*{Definição do Problema e Base de Dados}
O projeto atuará em um problema prático no qual a inferência em tempo real e de baixo consumo energético seja um limitante. Apresentam-se duas opções:

\vspace{0.5em}
\noindent \textbf{Opção 1: Detecção de ECG Fetal}
\begin{itemize}
    \item \textbf{Problema:} Extrair e classificar eletrocardiogramas fetais a partir de sinais mistos originários de registros abdominais maternos. É uma tarefa complexa de separação de sinal 1D afetada por ruídos de músculos e respiração.
    \item \textbf{Dataset Sugerido:} \textit{PhysioNet Abdominal and Direct Fetal ECG Database}.
    \item \textbf{Aplicação:} A série temporal é convertida em espectrogramas 2D, e a CNN atua como classificador de eventos. O referencial teórico já demonstrou alta eficácia, adequando-se perfeitamente a aceleradores embarcados.
\end{itemize}

\vspace{0.5em}
\noindent \textbf{Opção 2: Manutenção Preditiva Industrial}
\begin{itemize}
    \item \textbf{Problema:} Diagnóstico do estado de saúde de maquinário através da vibração dos componentes, identificando falhas mecânicas.
    \item \textbf{Dataset Sugerido:} \textit{CWRU Bearing Dataset} (Séries temporais de rolamentos).
    \item \textbf{Aplicação:} Transformação do sinal de vibração via STFT e classificação das anomalias usando a CNN com núcleos reduzidos, o que substitui a necessidade de modelos profundos 1D ou Redes Recorrentes (RNN/LSTM) pesadas.
\end{itemize}

\end{document}
